\section{Kiểm tra chất lượng dữ liệu và thiết kế đánh giá}
\label{sec:audit}

\subsection{Kiểm tra tính đầy đủ và tính hợp lệ}

Bước audit kiểm tra ba khía cạnh: tính liên tục của lịch, tính duy nhất của ngày và
miền giá trị của lưu lượng. Kết quả được tóm tắt trong
Bảng~\ref{tab:data-quality}. Chuỗi có đúng \DatasetN{} ngày, không thiếu ngày nào
trong khoảng \DateMin{}--\DateMax{}, không có ngày trùng và không có giá trị null.
Các giá trị lưu lượng đều hữu hạn, dương và không có ngày nào bằng 0. Điều này cho
thấy dữ liệu đã được làm sạch ở mức cơ bản và không cần nội suy hoặc loại bỏ quan
trắc.

\begin{table}[H]
\centering
\small
\begin{tabular}{@{}ll@{}}
\toprule
Hạng mục & Giá trị \\
\midrule
Số ngày quan trắc & 5113 \\
Phạm vi & 2010-01-01 → 2023-12-31 \\
Ngày thiếu & 0 \\
Ngày trùng & 0 \\
Giá trị null & 0 \\
Giá trị không hữu hạn & 0 \\
Giá trị âm & 0 \\
Giá trị bằng 0 & 0 \\
Chuỗi zero dài nhất & 0 \\
Trung bình (m³/s) & 295.6839 \\
Độ lệch chuẩn mẫu (m³/s) & 99.3151 \\
Nhỏ nhất (m³/s) & 94.3909 \\
Q95 (m³/s) & 417.230854 \\
Lớn nhất (m³/s) & 457.486 \\
SHA256 & f83d984747fe12e5\ldots{} \\
\bottomrule
\end{tabular}

\caption{Kết quả kiểm tra chất lượng dữ liệu gốc.}
\label{tab:data-quality}
\end{table}

Một điểm cần lưu ý là báo cáo chỉ kiểm tra chất lượng \emph{kỹ thuật} của tệp dữ liệu,
không kiểm tra chất lượng \emph{đo đạc} tại trạm. Việc không có ngày thiếu không có
nghĩa là số liệu đã được hiệu chỉnh theo cùng một quy trình trong suốt 14 năm. Vì
vậy các kết luận về xu hướng được trình bày thận trọng, đặc biệt khi so sánh giữa
các giai đoạn có mức nền rất khác nhau.

Mã băm SHA256 rút gọn của dữ liệu gốc là \texttt{f83d984747fe}; mã đầy đủ được lưu
trong \texttt{data\_quality.json} và được dùng để kiểm tra rằng
các lần chạy pipeline sau đọc đúng cùng một tệp. Nếu dữ liệu bị thay đổi dù chỉ một
ký tự, mã băm sẽ khác và báo cáo không còn tương ứng với kết quả cũ.

\subsection{Thiết kế chia tập và chống rò rỉ dữ liệu}

Bài toán được định nghĩa như sau: tại gốc thời gian $t$, mô hình nhận các quan sát
đến hết ngày $t$ và phát hành đồng thời bảy dự báo cho $t+1,\ldots,t+7$. Mỗi cặp
$(t,h)$ là một mẫu đánh giá. Cách định nghĩa này phản ánh đúng việc vận hành: tại
thời điểm ra quyết định, ta không biết giá trị tương lai, nhưng có thể dùng toàn bộ
lịch sử đã quan sát.

Các tập được chia theo thời gian, không xáo trộn:

\begin{table}[H]
\centering
\small
\begin{tabular}{llr}
\toprule
Tập & Khoảng ngày mục tiêu & Số gốc dự báo \\
\midrule
Train & 01/01/2010–31/12/2019 & 3.615 \\
Validation & 01/01/2020–31/12/2021 & 725 \\
Refit & 01/01/2010–31/12/2021 & 4.346 \\
Test & 01/01/2022–31/12/2023 & 724 \\
\bottomrule
\end{tabular}

\caption{Các tập huấn luyện, validation, refit và test.}
\label{tab:split}
\end{table}

Tập Train dùng để chọn tham số LightGBM. Tập Validation dùng để chọn cấu hình và
đánh giá sơ bộ trong quá trình phát triển. Tập Test chỉ được dùng một lần sau khi
cấu hình đã được khóa. Tập Refit gộp Train và Validation để huấn luyện mô hình cuối
cùng trước khi dự báo Test. Vì Test nằm hoàn toàn sau Refit, không có quan sát nào
của Test được dùng để tính lag, rolling window hoặc chọn tham số.

Số gốc dự báo nhỏ hơn số ngày của mỗi giai đoạn vì hai ràng buộc: cần tối thiểu 30
ngày lịch sử để tính đặc trưng, và cần đủ 7 ngày tương lai để tạo nhãn. Với Train,
gốc đầu tiên là 31/01/2010 và gốc cuối là 24/12/2019; với Validation, gốc đầu tiên
là 31/12/2019 và gốc cuối là 24/12/2021; với Test, gốc đầu tiên là 31/12/2021 và
gốc cuối là 24/12/2023. Các gốc này tạo ra 3.615, 725 và 724 mẫu tương ứng. Tổng
số cặp dự báo trên Test là $724 \times 7 = 5.068$ cho mỗi mô hình.

\subsection{Các kiểm tra chống rò rỉ}

Pipeline thực hiện bốn kiểm tra tự động:

\begin{itemize}[leftmargin=1.4em]
  \item \textbf{Không xáo trộn:} thứ tự thời gian được giữ nguyên; không có phép
  \texttt{train\_test\_\allowbreak split} ngẫu nhiên.
  \item \textbf{Biên tập:} ngày mục tiêu lớn nhất của Train không vượt quá
  31/12/2019; mục tiêu của Validation và Test nằm trong đúng khoảng đã khai báo.
  \item \textbf{Đặc trưng lịch:} các đặc trưng như tháng, ngày trong năm của
  \emph{ngày mục tiêu} được phép dùng vì chúng biết trước; không dùng bất kỳ giá trị
  lưu lượng nào của ngày mục tiêu.
  \item \textbf{Khóa cấu hình:} sau khi chọn mô hình, cấu hình và split được ghi
  vào manifest; các bước evaluate và export chỉ đọc lại cấu hình đã khóa.
\end{itemize}

Các kiểm tra này không bảo đảm mô hình tốt, nhưng bảo đảm kết quả Test là đánh giá
ngoài mẫu đúng nghĩa. Đây là yêu cầu tối thiểu để một báo cáo dự báo chuỗi thời gian
có thể được xem là nghiêm túc.

\subsection{Điểm mạnh và giới hạn của thiết kế}

Thiết kế theo thời gian có ưu điểm là mô phỏng đúng điều kiện triển khai, nhưng cũng
có giới hạn. Thứ nhất, chỉ có hai năm Test nên các chỉ số có thể dao động mạnh nếu
một năm có chế độ dòng chảy khác biệt. Thứ hai, vì dữ liệu chỉ có một chuỗi, không
thể đánh giá khả năng tổng quát hóa sang trạm khác. Thứ ba, các đặc trưng chưa bao
gồm biến khí tượng hoặc vận hành hồ chứa, nên mô hình chỉ học được cấu trúc nội tại
của chuỗi. Những giới hạn này được quay lại trong phần thảo luận.
